\documentclass[11pt,a4paper]{article}
\usepackage[margin=1in]{geometry}
\usepackage{amsmath,amssymb,mathtools}
\usepackage{bm}
\usepackage{microtype}
\usepackage{enumitem}

\newcommand{\R}{\mathbb{R}}
\newcommand{\E}{\mathbb{E}}
\newcommand{\softmax}{\mathrm{softmax}}
\newcommand{\diag}{\mathrm{diag}}

\title{PropertyScanner: Retrieval-Augmented Log-Residual Valuation with Cross-Attention}
\author{}
\date{}

\begin{document}
\maketitle

\begin{abstract}
This paper documents the valuation model and verification harness used in
PropertyScanner. The pipeline is local-first and market-anchored: retrieve
time-safe comparables, normalize by hedonic or registry indices, compute a
robust comp baseline, and predict log-residual quantiles via cross-attention.
Model outputs are blended with income signals and area intelligence, and
optionally calibrated with conformal prediction. The goal is to produce
interval-valued, auditable estimates suitable for analyst triage.
\end{abstract}

\section{Assumptions and Regimes}
\label{sec:assumptions}
We model residential property valuation under:
\begin{itemize}[leftmargin=1.5em]
  \item local-first execution with SQLite as system of record,
  \item sparse and drifting markets with limited sold labels,
  \item comp selection bounded by geo/structure filters and time-safe constraints,
  \item optional multimodal signals (text/image) that must degrade gracefully.
\end{itemize}
Assumptions validated by tests are called out in Section \ref{sec:verification}.
Assumptions sourced from literature are referenced in Section \ref{sec:method}
and grounded in hedonic and repeat-sales index literature
\cite{Rosen1974,BaileyMuthNourse1963,CaseShiller1987}.
Open assumptions are explicitly labeled in the verification log.

\section{Method}
\label{sec:method}
\subsection{Comparable Retrieval}
We retrieve $K$ comparables $\mathcal{C}$ for a target listing based on
semantic similarity and structural constraints. The search uses dense text
embeddings with strict filters for property type, size ratio, and time safety.

\subsection{Time Normalization}
Comparable prices are normalized to the valuation month using market indices:
\begin{equation}
\label{eq:time_adjust}
 y_j^{\ast} = y_j \cdot \frac{I_{r,\ell_0}(\tau_0)}{I_{r,\ell_j}(\tau_j)}.
\end{equation}
When indices are missing, the system falls back to raw prices and records the
fallback status in evidence metadata.

\subsection{Robust Comp Baseline}
We compute a robust baseline using a MAD-filtered weighted median:
\begin{equation}
\label{eq:baseline_mad}
 b = \mathrm{wmed}(\{y_j^{\ast}\}, \{\tilde{w}_j\})\quad\text{with}\quad
|y_j^{\ast} - m| \le 3\,\mathrm{MAD}.
\end{equation}
This baseline is the anchor for residual modeling and outlier resistance.

\subsection{Log-Residual Quantile Modeling}
Residuals are modeled in log space:
\begin{equation}
\label{eq:residual}
 \rho = \log y_0^{\ast} - \log b.
\end{equation}
Quantile predictions are reconstructed as:
\begin{equation}
\label{eq:price_quantile}
 \hat{y}_q = \exp(\log b + \hat{\rho}_q)\quad\text{for}\quad q \in \{0.1, 0.5, 0.9\}.
\end{equation}
Quantile modeling follows the regression-quantile framing of
\cite{KoenkerBassett1978}.
The fusion model uses cross-attention between target and comps to produce
residual quantiles \cite{Vaswani2017,ReimersGurevych2019}; strict baseline
fallbacks are applied on failure.

\subsection{Income Blend and Area Adjustment}
Income value is computed from rent and yield:
\begin{equation}
\label{eq:income_value}
 v_{\text{inc}} = \frac{12\,\hat{r}}{y_m / 100}.
\end{equation}
The final value is blended and adjusted by area intelligence with a bounded
adjustment factor $\delta$:
\begin{equation}
\label{eq:area_adjust}
 v' = v (1 + \delta), \quad u' = u (1 + |\delta|).
\end{equation}

\subsection{Uncertainty}
The reported uncertainty is the relative half-width:
\begin{equation}
\label{eq:uncertainty}
 u = \frac{\hat{y}_{0.9} - \hat{y}_{0.1}}{2\hat{y}_{0.5}}.
\end{equation}
Interval calibration uses asymmetric conformal prediction where available
\cite{RomanoPattersonCandes2019,BarberEtAl2021,AngelopoulosBates2021}.

\section{Implementation}
\label{sec:implementation}
We map equations to code entrypoints in
\texttt{paper/implementation\_map.md}. Key entrypoints include
\texttt{src/market/services/hedonic\_index.py} for Equation
\ref{eq:time_adjust}, \texttt{src/valuation/services/valuation.py} for
Equations \ref{eq:baseline_mad} and \ref{eq:price_quantile}, and
\texttt{src/valuation/services/conformal\_calibrator.py} for monotonicity
enforcement.

\section{Verification}
\label{sec:verification}
The verification harness provides:
\begin{itemize}[leftmargin=1.5em]
  \item sanity checks for time adjustment and interval ordering,
  \item invariants for monotonic quantiles and positive adjustment factors,
  \item regression tests pinned to \texttt{paper/artifacts/sanity\_case.json},
  \item a contract drift check for \texttt{paper/implementation\_map.md}.
\end{itemize}
Tests are implemented under \texttt{tests/unit/paper/} and executed via the
repo's pytest runner. Claims and their verification status are recorded in
\texttt{paper/verification\_log.md}.

\section{Reproducibility}
\label{sec:repro}
All reproducible commands are listed in \texttt{paper/README.md} and sourced
from repo scripts or CI usage patterns. The artifact generator produces a
small deterministic JSON sanity case used by regression tests.

\section{Limitations}
Current verification focuses on structural and interval properties, not full
end-to-end valuation accuracy. Coverage guarantees remain marginal unless
segment coverage checks are added.

\section{Conclusion}
This paper and harness provide a falsifiable, code-linked specification of the
PropertyScanner valuation pipeline. Future work should expand coverage metrics
and add ablation benchmarks.

\bibliographystyle{plain}
\bibliography{references}
\end{document}
